(1) 由参数化微分 \(dx/d\theta=a(1+\cos\theta)\)、\(dy/d\theta=a\sin\theta\)，得 \(ds/d\theta=\sqrt{x_\theta^2+y_\theta^2}=2a\cos(\theta/2)\)，因为 \(|\theta|<\pi\) 时余弦为正。积分并以 \(s(0)=0\) 定常数，得 \(s=4a\sin(\theta/2)\)；弧长符号与 \(\theta\) 相同，允许把曲线动力学化为单一广义坐标。
(2) 高度 \(y=a(1-\cos\theta)=2a\sin^2(\theta/2)=s^2/(8a)\)，因此 \(U(s)=mg s^2/(8a)\) 是严格的二次势而不需要小角展开。沿曲线的动能 \(T=m\dot s^2/2\)，Euler–Lagrange 方程为 \(m\ddot s+mg s/(4a)=0\)，即 \(\ddot s+\omega^2s=0\)，\(\omega^2=g/(4a)\)。
(3) 当振幅 \(|s_0|<4a\) 且不碰到尖点，通解 \(s(t)=A\cos(\omega t)+B\sin(\omega t)\)，完整周期为 \(P=2\pi/\omega=4\pi\sqrt{a/g}\)，与 \(A,B\) 无关。这是摆线等时性，不是任意重力曲线都具有的性质；普通圆弧单摆只有小摆角极限才近似等时。
(4) 若从 \(s(0)=s_0\)、\(\dot s(0)=0\) 释放，解为 \(s(t)=s_0\cos(\omega t)\)。首次到达最低点时 \(\omega t=\pi/2\)，故 \(t_{\rm bottom}=\pi/(2\omega)=\pi\sqrt{a/g}\)，与释放点无关；独立的能量积分 \(\dot s^2=\omega^2(s_0^2-s^2)\) 也给出同样的反三角函数结果。需要轨道光滑、无摩擦并且摆线尖点未被穿越。
